\documentclass[10pt,a4paper]{amsart}
\usepackage[margin=19mm]{geometry}
\usepackage{amsmath,amssymb,mathtools}
\usepackage[scr=rsfs,cal=euler]{mathalfa}
\usepackage{xfrac}
\usepackage[unicode,hidelinks,bookmarks=false]{hyperref}
\newcommand{\ie}{i.e.\ }
\newcommand{\cf}{cf.\ }
\newcommand{\resp}{respectively\ }
\newcommand{\Subsection}{\subsection}
\providecommand{\nobreakdash}{\nobreak\hbox{-}}
\setlength{\emergencystretch}{2em}
\multlinegap0pt
\newtheorem{proposition}{Proposition}
    \def\d{{\textnormal d}}
    \def\i{{\textnormal i}}
    \def\r{{\mathbb R}}
\title{Simple and accurate approximations to the Riemann zeta function}
\author{Alexey Kuznetsov}
\date{Source: arXiv:2503.09519v2 (CC BY 4.0)}
\begin{document}
\maketitle

\noindent\emph{Source and licence.} Simple and accurate approximations to the Riemann zeta function. Version arXiv:2503.09519v2 (CC BY 4.0), distributed by arXiv under the Creative Commons Attribution 4.0 International licence, http://creativecommons.org/licenses/by/4.0/, as recorded in the arXiv licence metadata for this exact version and shown on the abstract page https://arxiv.org/abs/2503.09519v2. Full licence notice: \texttt{licenses/arxiv-2503.09519-CC-BY-4.0.txt}.\par\smallskip

\noindent\textit{Editorial scope.} Counted unit: the article's proposition prop:P\_m (source0.tex lines 518--521). The article's complete original proof of it is delivered with it (source0.tex lines 522--547, one \texttt{proof} environment of 26 lines). No mathematics is written, repaired or re-derived: the statement and the proof are the article's own lines, in the article's own notation, and no retained line is substituted or re-flowed.  Retained uncounted as the source-local context the proof needs: lines 465--469.  Nothing was omitted from any retained span, so the package carries no omission record.  The retained lines cite 2 works, whose entries are transcribed verbatim from the article's own bibliography source in the e-print and delivered with short editorial labels recorded in the item's reference mapping.  No retained line contains a figure environment, an image-inclusion command or drawing code, so no image is delivered and the package declares no asset dependency.  Editorial formatting only, in addition: an AMS \texttt{amsart} preamble that restates the article's own declarations and macros used by the retained lines, namely the environments \texttt{proposition} on separate counters, together with the standard packages those lines need.  The retained environments print in this document as Proposition 1 in order of appearance, whereas the article prints the same items with the numbering of its own full document.  The complete native source of the article as distributed in the arXiv e-print for this exact version is bundled unchanged as \texttt{sources/174648/source0.tex}.
\begin{equation}\label{eqn:Mordell_integral}
H(y):=\int\limits_{\r} \frac{e^{-2\pi x^2+ 2 \pi  \theta  x y }}{\cosh(\pi \theta x)}   \d x=
\frac{1}{\cos(\pi y)} \Big[\sqrt{2} \cos(\pi y/2) e^{-\frac{\pi \i}{8}
(4y^2+1)} -e^{-\frac{\pi \i}{4}}\Big].
\end{equation}

\begin{proposition}\label{prop:P_m} Assume that there exist polynomials $\{P_n\}_{0\le n \le m}$ that are orthogonal with respect to the linear functional $L$. Then 
the polynomial $P_m$  must satisfy
$P_m(x)=-x^m P_m(1/x)$ for all $x\neq 0$.
\end{proposition} 

\begin{proof} We begin by stating a preliminary result about determinants. Let $B=\{b_{i,j}\}_{1 \le i,j \le n}$ be an $n\times n$ matrix. Following \cite{Golyshev2007}, we denote by $B^{\tau}$ the transpose of  $B$  with respect to the anti-diagonal. In other words, $B^{\tau}=\{\tilde b_{i,j}\}_{1 \le i,j \le n}$ where $\tilde b_{i,j}=b_{n+1-j,n+1-i}$ for $1\le i,j \le n$. 
It is true that
\begin{equation}\label{eqn:det_B_tau}
{\textnormal{det}}(B) = {\textnormal{det}}(B^{\tau}).
\end{equation} 
This follows from the factorization $B^{\tau}={\textrm{J}} B^T {\textrm{J}}$ (see \cite{Golyshev2007}), where ${\textrm{J}}$ is an $n\times n$ exchange matrix, defined via ${\textrm{J}}_{i,j}=1$ if $i+j=n+1$ and ${\textrm{J}}_{i,j}=0$ otherwise. 

Now we can prove Proposition \ref{prop:P_m}. Since $H(y)$ is an even function (see \eqref{eqn:Mordell_integral}),   the moments  satisfy $\mu_{p,k}=\mu_{p,4p+1-k}$ for $0\le k \le 4p+1$. Using this fact and the well-known representation of orthogonal polynomials as determinants (see formula (2.2) in \cite{MARCELLAN2001}), we obtain
\begin{equation}\label{eqn:P_m_determinant}
P_m(x)= C \times  {\textnormal{det}}\begin{bmatrix}
\mu_0 & \mu_1 & \mu_2 & \cdots & \mu_{2p-1} &  \mu_{2p} & \mu_{2p} \\
\mu_1 & \mu_2 & \mu_3 & \cdots & \mu_{2p} & \mu_{2p} & \mu_{2p-1} \\
\mu_2 & \mu_3 & \mu_4 & \cdots & \mu_{2p} & \mu_{2p-1} & \mu_{2p-2} \\
\cdots  & \cdots  & \cdots  &   \cdots & \cdots & \cdots  & \cdots  \\
\mu_{2p-1} & \mu_{2p} & \mu_{2p}  & \cdots &\mu_{3} &\mu_{2} & \mu_{1} \\
\mu_{2p} & \mu_{2p} & \mu_{2p-1}  & \cdots &\mu_{2} &\mu_{1} & \mu_{0} \\
1 & x & x^2 & \cdots & x^{2p-1} & x^{2p} & x^{2p+1} 
\end{bmatrix}
\end{equation}
for some constant $C$. In the above formula, we used the simplified notation $\mu_k=\mu_{p,k}$. 
For $k=0,1,\dots,m$ we denote by $A_k$ the submatrix obtained from the matrix in \eqref{eqn:P_m_determinant} by removing the last row and the $(k+1)$-st column.  Performing Laplace expansion along the last row, we obtain
$$
P_m(x)=C \sum\limits_{k=0}^{m} (-1)^k {\textnormal{det}}(A_k) x^k. 
$$
It is straightforward  to verify that $A_k=A_{m-k}^{\tau}$. The desired result follows from \eqref{eqn:det_B_tau}.   
\end{proof}

\begin{thebibliography}{99}
\bibitem[Golysh]{Golyshev2007} V.~V. Golyshev and J.~Stienstra. \newblock {Fuchsian equations of type DN}. \newblock {\em Communications in Number Theory and Physics}, 1(2):323--346, 2007.
\bibitem[MARCEL]{MARCELLAN2001} F.~Marcell\'an and R.~\'Alvarez-Nodarse. \newblock {On the ``Favard theorem” and its extensions}. \newblock {\em Journal of Computational and Applied Mathematics}, 127(1):231--254, 2001. \newblock \url{https://doi.org/10.1016/S0377-0427(00)00497-0}.
\end{thebibliography}
\end{document}
